\subsection{Method: The Sim-to-Real Gap \quad / \quad \textthai{วิธีการ: ช่องว่างระหว่างการจำลองและความเป็นจริง}}
\begin{paracol}{2}
Developing RL for physics often involves a "Sim-to-Real"\index{Sim-to-Real@ช่องว่างระหว่างการจำลองและความเป็นจริง (Sim-to-Real)} challenge. Training an agent on a real fusion reactor or an expensive robot is risky. Instead, we train in a high-fidelity synthetic simulation. However, if the simulation does not include realistic "noise" or "latency," the agent may fail when deployed on real hardware.

\switchcolumn

\begin{thai}
บทที่ 7 แสดงให้เห็นว่า AI ไม่ได้ทำหน้าที่แค่สังเกตการณ์ แต่ยังสามารถ "ลงมือเข้าควบคุม" ระบบฟิสิกส์ได้จริง การฝึกในเครื่องจำลอง (Simulation) และนำไปประยุกต์ใช้ในโลกจริง (Sim-to-Real) คืออนาคตของเครื่องเร่งอนุภาคและเตาปฏิกรณ์ฟิวชัน
\end{thai}
\end{paracol}
